\section{Cross-Object Backdoor Validity}
\label{sec:theorem}

We establish the theoretical foundation for applying standard causal
adjustment methods to object-centric process data. The key question is:
does the multi-object OCEL structure invalidate the backdoor criterion
when the treatment and outcome variables derive from events referencing
different object types?

\begin{theorem}[Cross-Object Backdoor Criterion]
\label{thm:backdoor}
Let $\mathcal{M} = (\mathcal{V}, \mathcal{D}, \mathcal{F}, \mathcal{U})$
be an Object-Centric Causal Process Model satisfying
Assumptions~\ref{ass:markov}--\ref{ass:sufficiency}. Let $T$ and $Y$ be
treatment and outcome variables in $\mathcal{V}$ that may be associated with
events referencing different object types in $\pi_{\mathit{omap}}$. Let
$\mathbf{Z} \subset \mathcal{V} \setminus \{T, Y\}$ be a set of variables
satisfying the \emph{backdoor criterion} relative to $(T, Y)$ in $\mathcal{D}$:
\begin{enumerate}[(i)]
    \item No variable in $\mathbf{Z}$ is a descendant of $T$ in $\mathcal{D}$;
    \item $\mathbf{Z}$ blocks every path between $T$ and $Y$ that has an
          arrow into $T$ (i.e., every back-door path from $T$ to $Y$).
\end{enumerate}

Then the causal effect $P(Y \mid do(T=t))$ is identifiable from the observed
distribution and given by the \emph{backdoor adjustment formula}:
\begin{equation}
\label{eq:backdoor}
P(Y = y \mid do(T=t)) \;=\; \sum_{\mathbf{z}}
P(Y = y \mid T = t,\, \mathbf{Z} = \mathbf{z}) \cdot P(\mathbf{Z} = \mathbf{z}),
\end{equation}
even when $T$ and $Y$ are derived from events that reference objects of
different types in $\pi_{\mathit{omap}}$, provided that the OCEL-to-variable
mapping $\pi_{\mathit{omap}}$ preserves the conditional independence structure
of $\mathcal{D}$.
\end{theorem}

\begin{proof}[Proof Sketch]
The classical backdoor adjustment theorem \citep[Theorem~3.3.2]{pearl2009causality}
holds whenever the backdoor criterion is satisfied in the causal DAG $\mathcal{D}$
under Assumptions~\ref{ass:markov}--\ref{ass:sufficiency}. The OCEL setting
introduces one additional complication: variables $T$ and $Y$ may be
\emph{derived} from events that reference objects of different types, and the
aggregation step (event log to scalar process variables) could in principle
introduce spurious dependencies or destroy genuine ones.

We show this does not invalidate the backdoor criterion under the following
argument.

\medskip
\noindent\textbf{Step 1: Aggregation preserves conditional independence structure.}
The OCEL-to-variable mapping $\pi_{\mathit{omap}}$ aggregates events per
object instance into scalar process variables at the case level (e.g., mean
lead time per order). Formally, each variable $V_i \in \mathcal{V}$ is a
measurable function $V_i = g_i(\{e \in E : \mathit{ot}(\pi_{\mathit{omap}}(e))
= \tau_i\})$ for some aggregation function $g_i$ and object type $\tau_i$.

Provided that $g_i$ is applied consistently (i.e., the same deterministic
function for all cases), the aggregation is a surjective many-to-one mapping
from the event-level data to the case-level variable. Under this consistency
condition, the joint distribution $P(\mathcal{V})$ induced by the OCEL data
retains all conditional independence relations that are encoded in the
underlying structural process that generated $\mathcal{L}$.

\medskip
\noindent\textbf{Step 2: The Markov condition holds for derived variables.}
Under Step~1, the derived variable DAG $\mathcal{D}$ over $\mathcal{V}$
satisfies the Markov condition (Assumption~\ref{ass:markov}) with respect to
$P(\mathcal{V})$. This is because the structural equations $\mathcal{F}$
determine the data-generating process at the event level, and consistent
aggregation propagates the conditional independence structure to the
case-level variables.

\medskip
\noindent\textbf{Step 3: Cross-object-type edges are not excluded.}
The backdoor criterion (conditions (i) and (ii) in the theorem statement) is
a property of the DAG $\mathcal{D}$ over $\mathcal{V}$, not of the object
types in $\pi_{\mathit{omap}}$. Since Step~1 ensures that $\mathcal{D}$
correctly encodes the conditional independence structure of $P(\mathcal{V})$,
the criterion can reference edges between variables associated with different
object types without any additional correction.

\medskip
\noindent\textbf{Conclusion.} Given Steps~1--3, the conditions of Pearl's
classical backdoor adjustment theorem are satisfied for the derived variable
DAG $\mathcal{D}$, regardless of whether the variables in the adjustment
set $\mathbf{Z}$ derive from same-type or different-type objects in
$\pi_{\mathit{omap}}$. Therefore Equation~\eqref{eq:backdoor} holds.
$\square$
\end{proof}

\begin{remark}[Scope of the Consistency Condition]
The aggregation consistency condition in Step~1 is satisfied by any
deterministic, case-level summary statistic (mean, max, count, duration).
It may fail for aggregations that are themselves causally downstream of
the treatment (e.g., aggregating post-treatment events), which would
introduce selection bias. This is a standard concern in causal inference
with time-varying treatments and is orthogonal to the OCEL multi-object
structure per se.

Full formalisation of the preservation condition and extension to partial
observability of object interactions (where some events in $E$ may not be
observed for all object instances) are left for future work.
\end{remark}

\begin{remark}[Practical Implication]
Theorem~\ref{thm:backdoor} justifies applying existing causal estimation
libraries (e.g., DoWhy \citep{sharma2020dowhy}) to OCEL-derived variable
sets without modification, provided the OCEL-to-variable mapping is
consistent in the sense of Step~1. We verify this empirically in our
experimental evaluation: the backdoor-adjusted estimate recovers the
planted causal coefficient within the expected statistical error across
all evaluation settings.
\end{remark}
